\documentclass[lualatex, paper={90mm, 160mm}]{jlreq}
\usepackage{tikz}
% luatexja は jlreq[lualatex] が自動的に読み込むため不要
% lltjext: 縦書きminipage (<t>オプション) のために必要
\usepackage{lltjext}
\usetikzlibrary{calc}

\begin{document}
\pagestyle{empty} 

\begin{tikzpicture}[remember picture, overlay]
  % 1. 背景の幾何学模様
  \begin{scope}[opacity=0.15, color=black!70]
    \foreach \x in {-5,...,5} {
      \foreach \y in {-9,...,9} {
        \draw[very thin] ($(current page.center) + (\x*10mm, \y*10mm)$) rectangle +(10mm, 10mm);
        \draw[very thin] ($(current page.center) + (\x*10mm, \y*10mm)$) -- +(10mm, 10mm);
      }
    }
  \end{scope}

  % 2. 表紙の二重枠線
  \draw[line width=1.2pt, color=black!80] ($(current page.south west) + (6mm, 6mm)$) rectangle ($(current page.north east) - (6mm, 6mm)$);
  \draw[line width=0.3pt, color=black!80] ($(current page.south west) + (7.5mm, 7.5mm)$) rectangle ($(current page.north east) - (7.5mm, 7.5mm)$);

  % 3. タイトル（縦書き）
  \node[anchor=north] at ($(current page.center) + (15mm, 50mm)$) {
    % 修正箇所2: lltjextの拡張機能を用いて、<t>オプションで縦書きBoxを生成
    \begin{minipage}<t>{15em}
      \begin{center}
        {\fontsize{28pt}{40pt}\selectfont 走れメロス\par}
      \end{center}
    \end{minipage}
  };

  % 4. 作者名（縦書き）
  \node[anchor=north] at ($(current page.center) + (-15mm, -10mm)$) {
    % 修正箇所3: 同様に縦書きBoxを使用
    \begin{minipage}<t>{10em}
      \begin{center}
        {\fontsize{16pt}{24pt}\selectfont 太宰 治\par}
      \end{center}
    \end{minipage}
  };
\end{tikzpicture}

\end{document}